\section{Discussion and Future Directions}
\label{sec:discussion}

This section addresses \textbf{RQ3} by identifying critical research directions emerging from our ecosystem analysis gaps and architectural paradigm shifts. We focus on directions with existing academic foundations rather than speculative best practices.

\subsection{Critical Research Directions}

Our analysis identifies priority research directions addressing fundamental gaps and persistent challenges revealed through ecosystem analysis. We focus on directions with existing academic foundations and near-term deployment potential.

\noindent\textbf{Verifiable Compliance Proofs.}
Zero-knowledge proofs and verifiable credentials offer cryptographic pathways for privacy-preserving compliance verification \cite{ben-sasson2017scalablezeroknowledge,yang2024researchidentitydata,zhang2024blockchainzeroknowledge}. Recent surveys and benchmarking studies \footnote{Polyhedra Network. Proof Arena: A Comprehensive Benchmark for Zero-Knowledge Proofs. 2024. \url{https://proofarena.org}}\footnote{National Institute of Standards and Technology (NIST). Privacy-Enhancing Cryptography: Zero-Knowledge Proofs. 2025. \url{https://csrc.nist.gov/projects/privacy-enhancing-cryptography/zero-knowledge-proof}} evaluate ZKP frameworks across dimensions including proof generation time, verification latency, proof size, memory requirements, trusted setup assumptions, and developer accessibility. These analyses reveal significant performance-practicality tradeoffs: SNARKs (e.g., Groth16, PLONK) offer succinct proofs and fast verification but require trusted setup ceremonies and impose substantial prover overhead; STARKs eliminate trusted setup at the cost of larger proof sizes and increased verification time; Bulletproofs provide transparent setup with logarithmic proof sizes but slower verification. For compliance use cases requiring frequent proof generation (e.g., transaction screening, attestation verification), prover performance becomes a critical practical constraint. Recursive proof composition, hardware acceleration, and specialized circuits show promise for improving feasibility. Regulatory acceptance remains uncertain, and guidance varies across jurisdictions.

\noindent\textbf{Semantic Transaction Understanding.}
Machine learning approaches for semantic DeFi transaction interpretation show promise, with recent work on language models and graph representation learning demonstrating improved fraud detection \cite{sun2025ethereumfrauddetection,azad2025machinelearningblockchain}. Smart contract analysis research advances technical foundations \cite{bhargavan2016formalverificationsmart,chen2024safecheckdetectingsmart,liao2024smartcontractvulnerability}, though production integration and generalization across diverse protocols remain limited.

\noindent\textbf{Behavior-Based Risk Assessment.}
Temporal graph neural networks and behavioral profiling techniques show improved detection of sophisticated money laundering patterns spanning transaction sequences \cite{albrecht2019usecryptocurrenciesmoney,cheng2023evolvepathtracer,cheng2024earlydetectionmalicious,arnold2024insightscaveatsmining,ouyang2024bitcoinmoneylaundering,li2025globallocalgraphattention,lin2024denseflowspottingcryptocurrency}. These approaches require substantial training data and computational resources while facing ongoing adversarial robustness challenges \cite{rodriguezvalencia2025systematicreviewartificial}.

\noindent\textbf{Privacy-Preserving Collaborative Analytics.}
Secure multi-party computation, federated learning, and private set intersection offer theoretical pathways for VASP collaboration without direct data sharing \cite{bansod2022challengesmakingblockchain}. However, computational costs, incentive misalignment, and limited production deployments may constrain near-term adoption.

\noindent\textbf{Intent-Aware Monitoring and Emerging KYI Concepts.}
Intent-centric architectures require monitoring at intent expression, solver routing, and execution layers rather than solely analyzing transaction outcomes \cite{mao2025knowyourintent}. A future-facing direction is the emergence of intent-level compliance concepts often described as "Know Your Intent (KYI)" in academic discourse, even though this label is not yet standardized in industry practice. Research challenges include formalizing intent representations, defining accountable parties across solvers and agents, and designing privacy-preserving intent screening that can be audited without exposing sensitive strategy details. Standardized interfaces and responsibility allocation frameworks represent priority development areas, but regulatory guidance and practical deployment examples remain nascent.

\noindent\textbf{Cross-Chain Compliance Infrastructure.}
Cross-chain attribution techniques including behavioral correlation and bridge monitoring advance technical foundations \cite{yousaf2019tracingtransactionscryptocurrency,zamyatin2021sokcommunicationdistributed,zhang2024securitycrosschainbridges,li2025blockchaincrosschainbridge,lin2024crosschainabnormaltransaction}, as detailed in \autoref{sec:ecosystem}. Standardization of data models and attribution databases faces coordination challenges. Modular blockchain architectures create additional monitoring complexity requiring execution layer access \cite{xu2024exploringblockchaintechnology,chaliasos2024analyzingbenchmarkingzkrollups}.

\noindent\textbf{Standardized Evaluation and Benchmarking.}
Standardized benchmarks for Web3 compliance remain limited, with research efforts employing incompatible datasets and evaluation metrics \cite{huda2025amlnet}. Ground truth labels for illicit blockchain activities are scarce, limiting rigorous model validation \cite{sankaewtong2025sokadvances}. Collaborative dataset curation and multi-dimensional evaluation frameworks capturing accuracy, false positive rates, latency, and adversarial robustness would advance research quality, though coordination challenges and privacy concerns limit progress.
